\renewcommand*{\authorOfNextLines}{Helmut Quinz}
Das indirekte oder sinngemäße Zitat gibt die Aussage(n) der genannten Quelle sinngemäß, jedoch mit eigenen Worten formuliert wieder.
Für das indirekte Zitat sind keine Anführungszeichen zu verwenden!
\subsubsection{Beispiel für ein direktes Zitat}
Betreffend häufig gehegter \gqq{Wahrheitsansprüche} befindet der Konstruktivismus, 
dass Weltbilder und damit wahrgenommene \gqqEmph{Wahrheiten} durch Interaktionen des Individuums mit seiner physischen und sozialen Umwelt erfahren und 
im Anschluss aktiv (wenn auch meist unbewusst) \gqq{konstruiert} werden. 
Erzählungen von der Wirklichkeit verlieren so gesehen ihre Absolutheit.
Sie werden funktionalistisch betrachtet und dadurch relativiert.(\cite[Seiten 68ff]{Simon2013})\\

Selbst Messungen bestätigen lediglich die Reproduzierbarkeit eines Vorgangs. 
Die Interpretation von Ursache und Wirkung bleibt eine eben solche.
Für die praktische Anwendung ist dies jedoch meist ausreichend!
\FloatBarrier